% Related work. Every characterisation here is drawn from the cited work's own abstract or
% published record, retrieved by DOI or from the arXiv API; see manuscript/tmlr/TODO_CITATIONS.md
% for which papers have been read in full and which have not.

\section{Related work}
\label{sec:related}

\paragraph{Whether a benchmark measures what it claims.}
One line of work treats benchmarks as instruments to be assessed rather than scoreboards to be
topped. BetterBench grades benchmarks against criteria for design, implementation and
documentation, and finds that many widely used ones score poorly \citep{reuel2024betterbench}.
A complementary framework asks which evaluative claims a given piece of evidence actually
licenses, distinguishing performance on a narrow test from the broad capability the
performance is often taken to demonstrate \citep{salaudeen2025measurement}. That work sits
upstream of this one. It asks whether the questions and the construct are right; we take the
benchmark and its reference answers as given and ask whether the ranking the pipeline
publishes is identified by the pipeline's own outputs. The concerns compound rather than
compete. A benchmark can be a valid instrument and still publish an unidentified ranking, and
a benchmark whose construct is shaky can still be scored unambiguously.

\paragraph{Configuration moves scores, and extraction is part of the configuration.}
That evaluation results move with choices nobody reports is documented for prompt formatting
\citep{sclar2023quantifying}, for benign perturbations that reorder leaderboards
\citep{alzahrani2024benchmarks}, and across the whole surface of a widely used evaluation
harness \citep{biderman2024lessons}. Answer extraction belongs to that surface. The pattern
the leaderboard's original scorer matches descends from the convention introduced with Minerva
\citep{lewkowycz2022minerva}, and proposals for better extractors, including learned ones,
follow \citep{yu2024xfinder}. The board studied here supplies its own demonstration: its MATH
results were rescored with a different verifier in February~2025, and the published scores
changed while the per-item verdict files did not \citep{kydlicek2025fixing}. Our claim is not
that a better extractor exists. It is that no extractor is neutral, that the choice is
currently unreported alongside the score, and that its consequence for a ranking can be
bounded rather than argued about. A better extractor changes the cell structure; it does not
remove the question of what happens where extractors disagree.

\paragraph{Verifiers as objects of measurement.}
Recent work measures verifiers directly, building benchmarks for verifier accuracy across
domains \citep{li2025verifybench} and training verifiers meant to be robust to the formats
that defeat regular expressions \citep{liu2025compassverifier}. That work and this one point
in opposite directions from the same observation. Given labels, one can rank verifiers and
adopt the best; we ask what survives when no verifier is adopted as truth, which is the
situation a reader of a published leaderboard is in. The two are complementary: a better
verifier shrinks the discordant cells, and the analysis here says what that shrinkage buys in
certified orderings.

\paragraph{Partial identification when the label is corrupted.}
The bounds in Section~\ref{sec:setup} are partial-identification statements in the sense of
\citet{manski2003partial}. The closest classical results concern exactly this structure: what
can be learned about a distribution when an unknown fraction of the data is contaminated
\citep{horowitz1995identification}, and the sharp bounds available when a binary variable is
misclassified at unknown rates \citep{molinari2008partial}. A verifier is a misclassifying
measurement of correctness whose error rate is unknown and may depend on the model, so the
same machinery applies with the models playing the role of sub-populations.

Concurrent work applies partial identification directly to language-model labels.
\citet{chen2026partial} study the prevalence of a latent binary attribute when the only
evidence is a panel of model reports whose errors may be arbitrarily dependent given the
truth, show that the prevalence is completely unidentified without restrictions that separate
the latent components, and obtain bounds under weak stochastic-ordering restrictions. The
diagnosis is the same as ours, and the resolutions differ in a way worth stating. They use
replication of the reports and a stochastic-ordering restriction, and target a single
prevalence with no ground truth at all. We use a human reference on a small audited subsample
and a bounded-heterogeneity restriction across models, and target differences of prevalences,
because a leaderboard reader consumes orderings rather than levels.

\paragraph{Reporting how far an assumption can be pushed.}
Assumption~\ref{ass:h} is not testable at this budget in any strong sense, so the useful thing
to report is how much of it a conclusion needs. That is the move made by design sensitivity in
observational studies, where a finding is described by the magnitude of unmeasured bias that
would be required to explain it away, and designs are compared by how much bias they can
tolerate \citep{rosenbaum2004design}. The breakdown value of
Definition~\ref{def:breakdown} is the same object in this setting, with cross-model
heterogeneity in place of unmeasured confounding: $\varepsilon^{*}_{jk}$ is the amount of
heterogeneity it would take to overturn the ordering $j \succ k$. The analogy also sets the
reporting standard we follow. A sensitivity parameter is worth little as a caveat in a
discussion section and worth a great deal as a number attached to each conclusion.

\paragraph{Two-phase sampling, verification bias, and raters without a gold standard.}
Auditing a cheap classifier with a small expensive sample is old. Two-phase sampling, in which
a large first-phase sample is measured cheaply and a small second phase is measured accurately,
goes back to \citet{neyman1938contribution}, and the binomial case with misclassification,
where a fallible classifier is calibrated against a subsample scored by an infallible one, is
\citeauthor{tenenbein1970double}'s double sampling scheme \citep{tenenbein1970double}.
The complication that the subsample is chosen by the fallible classifier's own verdict is
verification bias, and corrections for it are standard in diagnostic testing
\citep{begg1983assessment}; confidence limits that propagate uncertainty in the estimated
error rates into the corrected prevalence are also available \citep{lang2014confidence}.

Two features of our setting break the classical correction. The human reference is applied
to a sample stratified by the agreement pattern of several verifiers rather than by one
verifier's verdict, and the error rates are not assumed common across the populations being
compared, which is precisely the assumption Proposition~\ref{prop:oneverifier} shows to be
fatal to the question. The multi-rater literature is closer on the first point. Estimating
error rates from several imperfect tests without a gold standard is classical
\citep{hui1980estimating}, as is estimating observer error rates from repeated labels
\citep{dawid1979maximum}. Those methods identify rates by assuming the raters err
independently given the truth. We do not assume that, and Section~\ref{sec:data} reports what
the three verifiers share; instead of buying identification with an independence assumption,
we buy partial identification with agreement stratification and pay for the rest in interval
width.

\paragraph{Label-efficient evaluation with a proxy.}
Prediction-powered inference combines an abundant machine-generated proxy with a small
human-labeled sample to produce valid confidence intervals for a target parameter, using the
labeled sample to correct the proxy's bias \citep{angelopoulos2023ppi}. The family has grown
towards evaluation: stratified variants for hybrid evaluation of language models
\citep{fisch2024stratified}, active designs that spend the labeling budget where the proxy is
least certain \citep{zrnic2024active}, and corrections for evaluation against synthetic labels
\citep{boyeau2025autoeval}. Two papers target our object directly.
\citet{chatzi2024ranking} rank models from pairwise comparisons produced by a strong model,
calibrated against a small set of human comparisons, and return rank sets with a coverage
guarantee. \citet{lee2026judge} correct the bias a judge's imperfect sensitivity and
specificity induce in a naive score, build intervals that account for both the test set and
the human-labeled calibration set, and allocate calibration samples adaptively.

The structural relationship is close and the difference is the assumption, not the
arithmetic. Bias correction of this kind requires the proxy's operating characteristics,
estimated on the calibration sample, to transport to the population being scored. When those
characteristics are taken to be common across the models under comparison, which is the usual
calibration setup, Proposition~\ref{prop:oneverifier} applies: every pairwise gap is scaled by
one number, and a pair-specific reversal cannot occur. The reversals reported in
Section~\ref{sec:results} therefore sit outside that model. Our response is to stop correcting
and start bounding. Homogeneity is confined to the cells where independent verifiers agree,
it is carried as a parameter that is varied rather than a condition that is assumed, and the
output is a range with a breakdown value instead of a point estimate with an interval. The
price is honest and large: far fewer conclusions, each of them conditional on a number the
reader can disagree with.

\paragraph{Where to spend a label budget.}
The allocation question in Section~\ref{sec:data} is a stratified sampling problem, and the
classical answer optimises the variance of an estimated mean \citep{neyman1934stratified},
with unequal selection probabilities handled by the standard weighting
\citep{horvitz1952generalization} and the finite-population corrections used here taken from
the standard texts \citep{cochran1977sampling, kish1965survey}. Adaptive versions that learn
the within-stratum variances while sampling are current \citep{morikawa2026tsneyman}. Our
allocation is not a variance problem. The two uses of a label, narrowing an identified range
and testing the assumption on which every certified ordering depends, are not commensurable,
and the projection in Figure~\ref{fig:frontier} shows that the first is nearly flat in the
allocation. The budget went where variance minimisation would not have sent it.

\paragraph{What this paper adds.}
No estimator here is new, and the individual ingredients are all borrowed from the work above.
The contribution is the combination and the discipline: verifier provenance established from
the leaderboard's own records rather than inferred from agreement, agreement stratification
across $K$ verifiers so that pooling is confined to the cells where they agree, an
identification result that turns the pooling assumption into a reported parameter with a
breakdown value attached to every ordering, a design registered in advance that splits a small
label budget between certifying orderings and testing that assumption, and a released audit
with its quality gate and its failures.
